\section{完整名单、金融主体与二部图}
\subsection{名单是更新单位，不是任意十条记录}
机构—公司关系依赖名单口径。原始前十大股东、流通股东、有限售股东、发行前股东和表决权名单描述的范围不同，不能把同一报告期的记录直接叠加。主图优先使用InfoTypeCd=1且OrgDisSHFlg=1的原始总股东候选；流通2/2和补充记录另行保存。SHTypeCd描述股东类别，不作为总股东/流通名单分组键。

完整性不能只看排名。新源中同排名的后续比较单位有54,743个，其中54,601个持股数原值相同，这提示大量合法并列的可能性。按相同排名只保留第一行，会删掉真实不同股东；任取十行，又可能混入补录或不同版本。拟结合原名单标题、报告期、公开日、来源、SHSer和ShrSer\_Type规则识别整组，而不是在未知组键上进行去重。

\reportfigure{从原始总股东候选到接受版本与隔离记录}{fig:certification}{\FigCertification}{名单与身份认证流程设计；严格结构候选只用于定位，不自动标为完整组。}

图\ref{fig:certification}把组认证和主体认证分开。新包缺少更新时间和披露版本ID，需要通过原版、更正及数据库历史覆写说明补证。至少100个披露组与100个主体的检查应覆盖跨年、并列排名、缺ID、产品更名和名义托管，而不是只挑容易匹配的案例。检查结果按问题类型记录，不能只给一个“抽查通过”总数。

\subsection{主体粒度与历史证券对应}
GDID有助于跨期识别，但不能先验地视为稳定受益所有人。一个标识出现多个原名称，可能涉及更名、产品版本或名义账户；一个名称对应多个ID，也可能是不同产品。基金产品与管理人并非同一个投资节点，名义托管更不能直接代表共同决策。个人股东、名义托管和不能可靠确认的主体从机构图中排除，记录排除原因和覆盖变化。

\reporttable{主体及名单更新中的处理取舍}{tab:entityrules}{
\begin{tabularx}{\linewidth}{p{3cm}YY}\toprule
资料情况&拟采用处理&需要保留的证据\\\midrule
ID与名称跨期一致&作为主体候选继续核验&原名称、类别、公开资料对应\\
同ID多名称&不直接合并为稳定受益人&更名、产品、名义账户规则\\
缺ID的非自然人&待类型与名称跨期确认&原报告与唯一对应范围\\
个人或名义托管&不作为共同机构关系&类型/托管识别与排除原因\\
新完整名单&按生效日替换旧整个列表&完整组、原版/更正版本键\\
股东退出Top10&移除当前披露图边&不把该行为当作真实清仓\\\bottomrule
\end{tabularx}}

公司历史证券对应采用独立有效区间，不能由DRET观察起止推断上市和退市。未可靠连接的辅助记录不进入当天图，目标公司日仍保留并表达$M_g=0$。这种设计避免“有图的公司较容易预测”被样本筛选误计为图模型提升，年度及共同覆盖子样本统计则用于解释适用范围。

\section{比例、边龄和目标去自身的消息}
\subsection{边权与机构聚合}
对接受的完整版本，边记录公司$i$、主体$j$、总股本持股比例$p_{ij}$和市场日年龄$a_{ij}$。主口径使用有限、非负的总股本比例，不把其他百分号字段一律当作相同量纲。定义
\begin{equation}
 w_{ij}=p_{ij}\exp\left[-\operatorname{softplus}(\lambda_{edge})a_{ij}/252\right].
\end{equation}
公司序列表示经过可学习变换$\phi$进入机构总和：$S_j=\sum_kw_{kj}\phi(s_k)$，总权$d_j=\sum_kw_{kj}$。向目标$i$返回时扣除自身：
\begin{equation}
 u_{j\setminus i}=\frac{S_j-w_{ij}\phi(s_i)}{d_j-w_{ij}},\qquad d_j-w_{ij}>0.
\end{equation}
分母不存在有效其他公司时，该机构不返回关系证据。图\ref{fig:bipartite}突出的是返回路径的扣除，而不是从整个图网络中删掉目标自身。

\reportfigure{二部图中目标去自身的加权消息路径}{fig:bipartite}{\FigBipartite}{仅为公司与机构变量的算法示意，不是实际股东个案，也不是因果传播图。}

进一步将机构消息按目标边权汇总：
\begin{equation}
 r_i=\frac{\sum_{j\in\mathcal N(i)}w_{ij}\mathbf1_{d_j-w_{ij}>0}\psi(u_{j\setminus i})}
 {\sum_{j\in\mathcal N(i)}w_{ij}\mathbf1_{d_j-w_{ij}>0}},\qquad
 g_i=\tanh\{W_g[s_i;r_i]+b_g\}M_{g,i}.
\end{equation}
目标自项$s_i$与跨公司消息联合投影是合法的图编码；需要去掉的是同一$s_i$先进入机构、再绕路返回的重复信息。若没有正权共享邻居，$M_g=0$。图随机化和同编码器对照可以检查$g_i$的增量是否主要来自参数容量或目标自项。

\subsection{小邻域和数值边界}
去自身后只有一个其他公司时，机构消息退化为该公司的变换表示。其边权在机构内归一化中抵消，不能据此声称持股比例在每种拓扑中都影响消息。存在两个或更多其他公司时，比例和年龄改变其相对贡献；外层不同机构之间还保留目标边权的作用。表\ref{tab:peercases}有助于区分这几种情况。

\reporttable{去自身后不同邻域规模的含义}{tab:peercases}{
\begin{tabularx}{\linewidth}{p{3cm}YY}\toprule
机构邻域&返回消息&对关系证据的解释\\\midrule
只持有目标公司&分母为0，不返回&不生成跨公司信息\\
仅一个其他公司&$u=\phi(s_k)$&机构内权重抵消，外层机构选择仍有作用\\
两个其他公司&二者的加权凸组合&比例和不同边龄改变相对贡献\\
较大邻域&分段加权聚合&需报告采样/截断及全邻居校验差异\\\bottomrule
\end{tabularx}}

实现中采用有效权重和数值阈值保护除法，记录极小分母、非有限值和权重下溢影响的边数。现有原型使用$10^{-8}$级阈值；正式数据预检应与更高精度聚合抽查，不能把数值保护阈值解释成经济意义上的最低持仓要求。重复公司—主体边应在来源处理阶段拒绝，而不是在图层里悄悄相加。

先计算各机构总和，再对每条目标边减去自身，可以使去自身消息量级保持为$O(Ed)$。若对每个目标都重新遍历全部机构邻居，高度机构会产生大量重复计算。研究将按日合并目标与邻居，使用当前权重编码每个公司一次，再进行分段聚合。训练采样上限可从15/10起测，正式验证与推断优先使用同一规则的分块全邻居计算。

\section{归一化边龄为何不足以限制整体修正}
对同一完整名单，目标公司多条边通常共享其披露年龄$A_i$。若外层边权为$p_{ij}e^{-rA_i}$，归一化后
\begin{equation}
 \frac{p_{ij}e^{-rA_i}}{\sum_h p_{ih}e^{-rA_i}}=\frac{p_{ij}}{\sum_h p_{ih}}.
\end{equation}
共同年龄因子抵消。即使整个目标名单变旧，归一化外层消息未必自动变小。邻居公司各自不同年龄仍会影响机构内相对权重，但这与限制整体辅助影响是两个问题。

这一现象是设置图外年龄包络的直接理由。图内权重处理“哪些邻居更值得参考”，图外包络处理“整体关系资料最多能改变主干多少”。若只把边龄写进归一化均值，可能遗漏目标名单绝对陈旧度。相反，若全部只在图外衰减，又不能区分不同邻居资料的新旧。两级设计需要分别消融，不能把任何含年龄变量的网络都称为同一种时点机制。

\section{有界残差融合及可检查性质}
\subsection{将修正幅度和内容分开}
主干分数由$s_i$与同信息模型共同采用的市场、行业和诉讼背景$c_i$计算，$z_{s,i}=w_s^T[s_i;c_i]+b_s$。仅收益GRU仍不接收$c_i$。公告和图风险头输出$z_e,z_g$，内容修正为$\delta_m=\tanh(z_m)$。对$m\in\{e,g\}$，设置
\begin{align}
 a_m&=M_mQ_m e^{-\operatorname{softplus}(\lambda_m)A_m}
 \sigma\{G_m([s;h_m;M_m;A_m;Q_m])\},\\
 z&=z_s+a_e\delta_e+a_g\delta_g,\qquad p=\sigma(z).
\end{align}
三路向量维度128，标量门输入维度为$128+128+3=259$。公告年龄取最近有效事件，按60市场日归一化；图年龄取目标公司当前接受完整版本，按252日归一化。事件注意力和事件残差共享事件衰减参数，图边及图残差分别使用参数，共三个年龄率。

$M_m$表示有效辅助信息是否存在。$Q_e$来自格式和证据可追溯检查，不使用未来标签或模型自报置信度；主图不可靠边在入图前排除，$Q_g$取1。门的内容项允许辅助修正为正或负，因此“有界”不意味着所有公告只会提高风险。

\subsection{幅度上界与同模型回退}
由于$0\le\sigma(G_m)\le1$且$|\tanh(z_m)|\le1$，有
\begin{equation}
 |a_m\delta_m|\le M_mQ_m e^{-\operatorname{softplus}(\lambda_m)A_m}.
\end{equation}
全无有效公告和关系时，$M_e=M_g=0$，从而$z=z_s$。这是同一模型的回退；它不等于另行训练、没有共同结构控制或辅助损失的收益GRU。图\ref{fig:envelope}以一个明确给定的衰减率说明允许范围的变化。

\reportfigure{辅助logit修正的解析年龄包络}{fig:envelope}{\FigEnvelope}{纯数学示意：$M=Q=1$且衰减率$r=1$，曲线为$\pm e^{-A}$，不是从训练资料拟合出的参数，也不是预测性能曲线。}

\reporttable{解析包络的数值及其解释}{tab:envelopetable}{
\begin{tabularx}{\linewidth}{p{2.3cm}p{3cm}Y}\toprule
归一化年龄 $A$&$e^{-A}$（示意$r=1$）&可以说明什么\\\midrule
0&1.0000&新鲜资料允许的最大修正\\
0.5&0.6065&幅度包络缩小，不确定实际方向\\
1&0.3679&不能推出概率必然下降\\
2&0.1353&不能当作已训练模型的真实贡献\\\bottomrule
\end{tabularx}}

门本身依赖内容和年龄，实际修正可能小于包络，也可能因内容变化而改变符号。因此，年龄上界单调不代表风险概率单调。由sigmoid导数上界$1/4$和均值定理，还可得
\begin{equation}
 |p-\sigma(z_s)|\le\tfrac14(|a_e\delta_e|+|a_g\delta_g|).
\end{equation}
该界描述同模型允许的概率偏移，不涉及泛化误差，也不保证限制后的排序更好。

\subsection{约束的代价和可检验问题}
有界修正会改变模型能表达的函数。如果某份旧资料仍有强信息，过度衰减可能损害预测；若事件质量标记只反映接口格式，还可能错误排除有用正文。这些代价不能由上述不等式排除。衰减率采用可学习参数，可以接近很小值，但能否学到合理权重仍取决于训练数据和优化。

研究将比较普通门控、年龄仅作特征和包络约束，检查自然年龄分组及人为增龄压力。压力实验固定资料内容和模型权重，只改变输入年龄，观察上界、分数偏移和排序变化；自然分组则保留真实资料状态，两者回答不同问题。普通门控也可以加入相似约束，所以方法区别最终仍需同编码器、同信息和合理预算下的结果支持。

\section{联合损失、训练掩码和优化}
风险任务主损失为训练目标公司日的未加权BCE。有效支路另加辅助损失：
\begin{equation}
 \mathcal L=\mathcal L_{main}+0.2\mathcal L_s+0.1\mathcal L_e+0.1\mathcal L_g.
\end{equation}
每个辅助BCE只在当前训练目标与该支路可用的交集上平均，没有交集时该项为0。仅收益模型只用主BCE；去辅助损失控制保留融合结构，移除额外支路监督。这样不会把邻居未来标签当成可共享训练目标。

\reportfigure{共享日批次的当前权重编码、消息计算与监督路径}{fig:training}{\FigTraining}{联合训练设计；反向传播可经过邻居历史表示，但不额外监督邻居未来标签。}

Qwen先完成抽取评价再冻结；正文原始向量可缓存，事件投影、GRU、图消息与门在风险任务中共同学习。冻结随机GRU或随机图投影会让输入表示与任务不匹配，不能以“节省训练”代替正式联合学习。训练采用AdamW，学习率候选$3\times10^{-4}$和$10^{-3}$，weight decay为$10^{-4}$，最多20轮，验证AP连续3轮不改善早停，梯度范数上限1。

概率校准只在规定验证集拟合温度或截距，检查Brier及校准情况\cite{r30}。在校准与稳定性评价之前，输出宜称风险分数。所有配置、预处理参数、缓存键和随机种子与运行记录一起保存，便于追踪模型选取及恢复后的结果。
